\documentclass[a4paper,12pt]{article}
\usepackage[T2A]{fontenc}
\usepackage[utf8]{inputenc}
\usepackage[russian]{babel}
\usepackage{amsmath,amssymb,mathtools}
\renewcommand{\familydefault}{\sfdefault}
\usepackage{icomma}
\usepackage[a4paper,left=22mm,right=22mm,top=20mm,bottom=23mm]{geometry}
\usepackage{microtype}
\usepackage{tikz}
\usetikzlibrary{calc,arrows.meta,positioning,shapes.geometric}
\usepackage{pgfplots}
\pgfplotsset{compat=1.18}
\usepackage{tabularx,booktabs}
\usepackage{enumitem,needspace}
\usepackage{fancyhdr,setspace,xcolor}
\definecolor{accent}{HTML}{246C70}
\definecolor{darkaccent}{HTML}{16464C}
\definecolor{pale}{HTML}{DAE7E6}
\definecolor{textgray}{HTML}{333A3A}
\color{textgray}
\onehalfspacing
\setlength{\parindent}{0pt}
\setlength{\parskip}{4pt}
\setlist[enumerate]{leftmargin=*,itemsep=7pt,topsep=4pt}
\setlist[itemize]{leftmargin=*,itemsep=3pt,topsep=3pt}
\pagestyle{fancy}
\fancyhf{}
\renewcommand{\headrulewidth}{0pt}
\renewcommand{\footrulewidth}{0.25pt}
\fancyfoot[L]{\scriptsize \textcolor{darkaccent}{АННА БАННОВА \; / \; 11.10.26}}
\fancyfoot[R]{\scriptsize\thepage}
\setlength{\footskip}{13mm}
\newcommand{\topic}[1]{\Needspace{5\baselineskip}\vspace{4mm}{\Large\bfseries\color{darkaccent}#1}\par\vspace{1mm}\textcolor{pale}{\rule{\linewidth}{0.7pt}}\vspace{1mm}}
\newcommand{\subtopic}[1]{\Needspace{3\baselineskip}\vspace{2mm}{\normalsize\bfseries\color{accent}#1}\par\vspace{1mm}}
\newcommand{\important}[1]{\textbf{\color{darkaccent}#1}}
% Компактная координатная сетка тренировочной задачи.
% Аргументы: выражение в синтаксисе pgfplots, точка A, точка B,
% границы x_min, x_max, y_min, y_max.
\newcommand{\taskplot}[7]{%
\begin{tikzpicture}
\begin{axis}[
 width=\linewidth,height=53mm,
 xmin=#4,xmax=#5,ymin=#6,ymax=#7,
 axis lines=middle,axis line style={gray!70,line width=.45pt},
 grid=major,grid style={gray!24,line width=.3pt},
 xtick distance=1,ytick distance=1,
 tick label style={font=\scriptsize},
 label style={font=\scriptsize},
 xlabel={$x$},ylabel={$y$},
 xlabel style={at={(axis description cs:1.02,0.13)},anchor=west},
 ylabel style={at={(axis description cs:0.10,1.01)},anchor=south},
 clip=false,domain=#4:#5,samples=110,
 enlargelimits=false]
 \addplot[accent,line width=1.2pt] {#1};
 \addplot+[only marks,mark=*,mark size=2.2pt,color=darkaccent] coordinates {(#2) (#3)};
 \node[font=\scriptsize,fill=white,inner sep=1pt,anchor=south east] at (axis cs:#2) {$A$};
 \node[font=\scriptsize,fill=white,inner sep=1pt,anchor=south east] at (axis cs:#3) {$B$};
\end{axis}
\end{tikzpicture}}
\newcommand{\exercisegraph}[3]{%
\begin{minipage}[t]{.475\linewidth}
\raggedright\textbf{#1.}\quad #2\par\vspace{1mm}
#3
\end{minipage}}
\begin{document}
\thispagestyle{empty}
\begin{center}
\vspace*{34mm}
{\large\color{accent}\textbf{МАТЕМАТИКА \enspace / \enspace ЕГЭ}}\\[24mm]
{\Huge\bfseries\color{darkaccent}Чек-лист}\\[8mm]
{\LARGE\bfseries Показательная функция}\\[3mm]
{\Large График, узловые точки и поиск коэффициентов}\\[26mm]
{\normalsize\color{textgray}11 октября 2026 года}\\[8mm]
{\small Путь А \enspace\textcolor{accent}{\textbullet}\enspace 10 заданий}\\[33mm]
{\normalsize\bfseries Лёвин Артём Александрович}\\[3mm]
{\normalsize эксклюзивно для Анны Банновой}
\end{center}
\begin{tikzpicture}[remember picture,overlay]
\draw[accent!60,line width=1pt]
 ($(current page.south west)+(12mm,11mm)$)
 .. controls ($(current page.south west)+(75mm,35mm)$)
 and ($(current page.south east)+(-75mm,5mm)$)
 .. ($(current page.south east)+(-12mm,9mm)$);
\draw[pale,line width=.45pt]
 ($(current page.south west)+(12mm,5mm)$)
 .. controls ($(current page.south west)+(83mm,25mm)$)
 and ($(current page.south east)+(-73mm,5mm)$)
 .. ($(current page.south east)+(-12mm,5mm)$);
\end{tikzpicture}
\clearpage

\topic{1. Как прочитать условие}
\important{Функция} $y=f(x)$ сопоставляет каждому допустимому $x$ определённое значение $y$.
Запись $f(6)$ означает: вместо $x$ подставить $6$ и вычислить результат.
Если требуется найти $x$, при котором $f(x)=29$, число $29$ подставляют \emph{вместо значения функции}.

\subtopic{Два вида показательных функций}
В сегодняшних заданиях встречаются:
\[
\boxed{f(x)=a^x+b}\qquad\text{и}\qquad
\boxed{f(x)=a^{x+b}},\qquad a>0,\ a\ne1.
\]
\begin{center}
\renewcommand{\arraystretch}{1.27}
\begin{tabularx}{\linewidth}{@{}p{39mm}X X@{}}
\toprule
 & $a^x+b$ & $a^{x+b}$\\
\midrule
Где стоит $b$? & За пределами степени & В показателе степени\\
Удобная точка & $x=0$: $f(0)=1+b$ & $f(x)=1$: тогда $x+b=0$\\
Возможен ли $f(x)\le0$? & Да, при некоторых $b$ & Нет: всегда $f(x)>0$\\
\bottomrule
\end{tabularx}
\end{center}
\important{Скобки решают всё:} $a^x+b$ и $a^{x+b}$ --- разные функции.

\topic{2. Метод узловых точек}
\important{Узловая точка} --- точка графика с двумя целочисленными координатами: $A(x_A;y_A)$.
Первая координата --- $x$, вторая --- $y=f(x)$.

Если в формуле неизвестны два коэффициента $a$ и $b$, обычно достаточно двух подходящих точек графика.
Удобнее начинать с точек, где $x=0$, $x=1$ или $f(x)=1$: так проще вычислять степени.

\begin{enumerate}[label=\textbf{\arabic*.}]
\item Найдите \important{две точки, лежащие именно на графике}, и запишите их координаты.
\item Подставьте каждую пару $(x;y)$ в исходную формулу. Получится система из двух уравнений.
\item Используйте $a^0=1$, $a^1=a$, свойства отрицательных и дробных степеней, найдите $a$ и $b$.
\item Запишите окончательную функцию с найденными коэффициентами.
\item Выполните последнее требование: вычислите $f(c)$ или решите $f(x)=m$; проверьте $a>0$, $a\ne1$.
\end{enumerate}

\clearpage
\topic{3. Пример: $f(x)=a^x+b$}
На графике отмечены узловые точки $A(0;-2)$ и $B(1;-1)$.
Требуется найти $f(6)$.
\begin{center}
\begin{tikzpicture}
\begin{axis}[
 width=133mm,height=76mm,
 xmin=-2.5,xmax=3.5,ymin=-3.6,ymax=6.5,
 axis lines=middle,axis line style={gray!70},
 grid=major,grid style={gray!24},
 xtick distance=1,ytick distance=1,
 xlabel={$x$},ylabel={$y$},
 xlabel style={at={(axis description cs:1.02,.32)},anchor=west},
 ylabel style={at={(axis description cs:.31,1.01)},anchor=south},
 clip=false,domain=-2.3:3.25,samples=150]
 \addplot[accent,line width=1.35pt] {pow(2,x)-3};
 \addplot+[only marks,mark=*,mark size=2.1pt,color=darkaccent] coordinates {(0,-2) (1,-1) (2,1) (3,5)};
 \node[anchor=east,font=\small,fill=white,inner sep=1pt] at (axis cs:0,-2) {$A$};
 \node[anchor=north west,font=\small,fill=white,inner sep=1pt] at (axis cs:1,-1) {$B$};
 \node[anchor=north west,font=\small,fill=white,inner sep=1pt] at (axis cs:2,1) {$C$};
 \node[anchor=east,font=\small,fill=white,inner sep=1pt] at (axis cs:3,5) {$D$};
\end{axis}
\end{tikzpicture}
\end{center}
\subtopic{Решение через систему}
\[
\begin{cases}
-2=a^0+b,\\
-1=a^1+b
\end{cases}
\ \Longrightarrow\ 
\begin{cases}
-2=1+b,\\
-1=a+b
\end{cases}
\ \Longrightarrow\ 
\begin{cases}
b=-3,\\a=2.
\end{cases}
\]
Следовательно, $f(x)=2^x-3$ и
\[
\boxed{f(6)=2^6-3=64-3=61}.
\]
\important{Проверка:} $f(0)=1-3=-2$, $f(1)=2-3=-1$ --- обе исходные точки принадлежат найденному графику.

\topic{4. Основание меньше единицы}
Для $f(x)=a^x+b$ даны точки $A(-2;0)$ и $B(0;-3)$.
Из второй точки $-3=1+b$, поэтому $b=-4$. Первая точка даёт
\[
0=a^{-2}-4\ \Longrightarrow\ a^{-2}=4
\ \Longrightarrow\ a=4^{-1/2}=\frac12,
\]
где выбран \important{положительный} корень, поскольку $a>0$.
Получили $f(x)=\left(\frac12\right)^x-4=2^{-x}-4$.

\clearpage
\topic{5. Когда $b$ внутри показателя}
Дана функция $f(x)=a^{x+b}$; на графике точки $A(-3;1)$ и $B(1;4)$.
\[
\begin{cases}
1=a^{-3+b},\\
4=a^{1+b}.
\end{cases}
\]
Поскольку $a>0$ и $a\ne1$, равенство $a^t=1$ означает $t=0$.
Поэтому $-3+b=0$, откуда $b=3$.
\[
4=a^{1+3}=a^4
\ \Longrightarrow\
a=4^{1/4}=2^{1/2}.
\]
\[
\boxed{f(x)=\left(2^{1/2}\right)^{x+3}=2^{(x+3)/2}}.
\]
Если теперь нужно найти $x$, при котором $f(x)=8$, то
\[
2^{(x+3)/2}=2^3
\ \Longrightarrow\ \frac{x+3}{2}=3
\ \Longrightarrow\ \boxed{x=3}.
\]

\topic{6. Что повторить из предыдущего занятия}
\begin{center}
\renewcommand{\arraystretch}{1.25}
\begin{tabularx}{\linewidth}{@{}X X@{}}
\toprule
\textbf{Правило} & \textbf{Пример}\\
\midrule
Нулевая степень: $a^0=1$ для $a\ne0$ & $7^0=1$\\
Отрицательная степень: $a^{-r}=1/a^r$ & $2^{-3}=1/8$\\
Дробная степень: $a^{1/n}=\sqrt[n]{a}$ для $a>0$ & $9^{1/2}=3$\\
Степень степени: $(a^r)^s=a^{rs}$ для $a>0$ & $(4^{1/2})^2=4$\\
Равенство степеней: $a^u=a^v\Rightarrow u=v$ для $a>0$, $a\ne1$ & $2^x=32\Rightarrow x=5$\\
\bottomrule
\end{tabularx}
\end{center}
\subtopic{Типичные ошибки}
\begin{itemize}
\item Переставлять местами $x$ и $y$ при подстановке координат точки.
\item Считать, что $f(6)$ означает умножение $f$ на $6$.
\item Путать $a^x+b$ и $a^{x+b}$.
\item При решении $a^{2}=c$ забывать ограничение $a>0$ для основания показательной функции.
\item Удалять основания в уравнении $a^u=a^v$, не проверив $a>0$ и $a\ne1$.
\end{itemize}

\clearpage
\thispagestyle{plain}
\begin{center}
{\Large\bfseries\color{darkaccent}Алгоритм: восстановление функции по графику}\\[3mm]
{\small\color{accent}От узловых точек к ответу}
\end{center}
\vspace{0mm}
\begin{center}
\begin{tikzpicture}[
 >=Stealth,
 node distance=8mm,
 term/.style={ellipse,draw=accent,line width=.8pt,fill=pale!22,align=center,text width=75mm,minimum height=13mm,inner sep=3mm,font=\small},
 action/.style={rectangle,rounded corners=2pt,draw=accent!85,line width=.75pt,align=center,text width=86mm,minimum height=14mm,inner sep=3mm,font=\small},
 decision/.style={diamond,aspect=2.4,draw=darkaccent,line width=.8pt,align=center,text width=34mm,inner sep=1.5mm,font=\small},
 side/.style={rectangle,rounded corners=2pt,draw=accent!85,align=center,text width=54mm,minimum height=16mm,inner sep=3mm,font=\small},
 arr/.style={->,line width=.9pt,draw=accent}
]
\node[term] (start) {Даны график и формула\\ с неизвестными $a$ и $b$};
\node[action,below=of start] (dots) {Выберите две узловые точки\\ $A(x_1;y_1)$ и $B(x_2;y_2)$};
\node[decision,below=10mm of dots] (kind) {Как записана функция?};
\node[side,below left=12mm and 1mm of kind] (left) {$f(x)=a^x+b$\\ $y_i=a^{x_i}+b$};
\node[side,below right=12mm and 1mm of kind] (right) {$f(x)=a^{x+b}$\\ $y_i=a^{x_i+b}$};
\node[action,below=39mm of kind] (system) {Составьте систему и найдите $a$ и $b$.\\ Запишите формулу $f(x)$};
\node[action,below=11mm of system] (goal) {Нужно найти $f(c)$? Подставьте $x=c$.\\ Нужно найти $x$? Решите $f(x)=m$.};
\node[term,below=11mm of goal] (last) {Проверьте $a>0$, $a\ne1$,\\ координаты точек и ответ};
\draw[arr] (start)--(dots);
\draw[arr] (dots)--(kind);
\draw[arr] (kind.south west) -| (left.north);
\draw[arr] (kind.south east) -| (right.north);
\draw[arr] (left.south) |- (system.west);
\draw[arr] (right.south) |- (system.east);
\draw[arr] (system)--(goal);
\draw[arr] (goal)--(last);
\end{tikzpicture}
\end{center}
\clearpage

\begin{center}
{\LARGE\bfseries\color{darkaccent}Тренировочный блок}\\[2mm]
{\small Путь А \quad\textcolor{accent}{\textbullet}\quad 10 заданий}
\end{center}
\textbf{Порядок работы.} Для каждой задачи на график сначала выпишите координаты точек $A$ и $B$, затем найдите $a$ и $b$. Запишите вычисления и окончательный ответ. Масштаб делений указан на осях.
\begin{enumerate}[label=\textbf{\arabic*.}]
\item Найдите значение выражения:
\[16^{1/2}+8^{2/3}-2^{-2}.\]
\item Решите показательное уравнение:
\[4^{x-1}=8^{x/2}.\]
\end{enumerate}
\vspace{3mm}
\begin{center}
{\normalsize\bfseries\color{accent}Задания 3--6. По графику восстановите $f(x)=a^x+b$.}
\end{center}
\vspace{2mm}
\exercisegraph{3}{По графику $f(x)=a^x+b$ найдите $f(2)$.}{%
\taskplot{pow(3,x)-3}{0,-2}{1,0}{-1.7}{2.3}{-3.5}{7.5}}
\hfill
\exercisegraph{4}{По графику $f(x)=a^x+b$ найдите $f(3)$.}{%
\taskplot{pow(3,x)}{0,1}{1,3}{-1.7}{2.3}{-0.8}{9.5}}
\par\vspace{6mm}
\exercisegraph{5}{По графику $f(x)=a^x+b$ найдите $f(2)$.}{%
\taskplot{pow(0.5,x)-2}{-1,0}{0,-1}{-2.2}{2.8}{-2.5}{3.4}}
\hfill
\exercisegraph{6}{По графику $f(x)=a^x+b$ найдите $f(4)$.}{%
\taskplot{pow(10,x/2)-3}{0,-2}{2,7}{-1.5}{2.18}{-3.6}{8.5}}
\par
\clearpage
\begin{center}
{\LARGE\bfseries\color{darkaccent}Тренировка: продолжение}\\[2mm]
{\small Задания 7--10}
\end{center}
\vspace{2mm}
\exercisegraph{7}{Для $f(x)=a^x+b$ найдите $x$, при котором $f(x)=22$.}{%
\taskplot{pow(3,x)-5}{0,-4}{1,-2}{-1.5}{1.5}{-5.5}{1.4}}
\hfill
\exercisegraph{8}{Для $f(x)=a^{x+b}$ найдите $f(1)$.}{%
\taskplot{pow(2,x+2)}{-2,1}{0,4}{-3.3}{0.8}{-0.3}{8.2}}
\par\vspace{7mm}
\exercisegraph{9}{Для $f(x)=a^{x+b}$ найдите $x$, при котором $f(x)=32$.}{%
\taskplot{pow(2,x+3)}{-3,1}{0,8}{-4.3}{0.15}{-0.2}{9.7}}
\hfill
\exercisegraph{10}{Для $f(x)=a^{x+b}$ найдите $x$, при котором $f(x)=27$.}{%
\taskplot{pow(3,(x+1)/2)}{-1,1}{3,9}{-2.5}{3.18}{-0.3}{10.5}}
\par
\vfill
\important{Самопроверка хода решения:} в заданиях 7, 9 и 10 требуется найти аргумент $x$ по известному значению функции. После получения формулы решите показательное уравнение.
\clearpage
\begin{center}
{\LARGE\bfseries\color{darkaccent}Ответы для самопроверки}
\end{center}
\vspace{6mm}
\renewcommand{\arraystretch}{1.65}
\begin{center}
\begin{tabularx}{\linewidth}{@{}>{\bfseries}p{16mm}X@{}}
\toprule
№ & Ответ\\
\midrule
1 & $\dfrac{31}{4}$\\
2 & $x=4$\\
3 & $f(2)=6$\\
4 & $f(3)=27$\\
5 & $f(2)=-\dfrac74$\\
6 & $f(4)=97$\\
7 & $x=3$\\
8 & $f(1)=8$\\
9 & $x=2$\\
10 & $x=5$\\
\bottomrule
\end{tabularx}
\end{center}
\vspace{8mm}
\important{Проверьте дополнительно:} восстановленная функция проходит через обе заданные точки; основание положительно и отличается от единицы; в конце вычислена именно та величина, о которой спрашивали в задаче.
\vfill
\begin{center}
\small\color{darkaccent}Лёвин Артём Александрович эксклюзивно для Анны Банновой
\end{center}
\end{document}
